\chapter{Analiza problemu i cel projektu}
\label{rozdzial1}

Wspólne mieszkanie, czy to w gronie rodziny, czy ze współlokatorami, wiąże się z koniecznością koordynacji wielu aspektów życia codziennego. Rozdział ten poświęcony jest analizie problemów związanych z zarządzaniem gospodarstwem domowym, zdefiniowaniu celu projektowanej aplikacji oraz przeglądowi istniejących na rynku rozwiązań.

\section{Opis problemu współdzielenia obowiązków domowych}

Zarządzanie domem to złożony proces, na który składa się szereg powtarzalnych czynności oraz decyzji podejmowanych przez grupę osób. Do najczęstszych problemów, z jakimi borykają się domownicy, należą:

\begin{itemize}
    \item \textbf{Niejasny podział obowiązków:} Brak precyzyjnego przypisania zadań (np. sprzątanie łazienki, wynoszenie śmieci) prowadzi do sytuacji, w których obowiązki są zaniedbywane lub wykonywane ciągle przez te same osoby, co rodzi frustrację i konflikty.
    \item \textbf{Problemy z rozliczeniami finansowymi:} Wspólne wydatki na czynsz, media czy zakupy spożywcze wymagają sprawiedliwego podziału kosztów. Ręczne spisywanie wydatków jest uciążliwe i podatne na błędy, a brak terminowych spłat między współlokatorami zaburza płynność finansową.
    \item \textbf{Dublowanie zakupów lub braki w zapasach:} Brak zsynchronizowanej listy zakupów skutkuje kupowaniem niepotrzebnych produktów lub - przeciwnie - zapominaniem o artykułach pierwszej potrzeby.
    \item \textbf{Trudności w komunikacji i podejmowaniu decyzji:} Ustalenia dotyczące wspólnych spraw (np. wybór dostawcy internetu, planowanie posiłków) często giną w natłoku wiadomości na komunikatorach, co utrudnia osiągnięcie konsensusu.
\end{itemize}

Brak zorganizowanego systemu wsparcia tych procesów obciąża domowników tzw. ,,obciążeniem mentalnym'', polegającym na konieczności ciągłego pamiętania o terminach i zobowiązaniach.

\section{Cel i zakres aplikacji}

Głównym celem projektu jest stworzenie aplikacji mobilnej ,,Household Manager'', która zintegruje kluczowe funkcjonalności niezbędne do sprawnego zarządzania domem w jednym, intuicyjnym interfejsie. System ma za zadanie zminimalizować ryzyko konfliktów poprzez zapewnienie transparentności działań i automatyzację powtarzalnych czynności.

Zakres funkcjonalny aplikacji obejmuje:
\begin{enumerate}
    \item \textbf{Zarządzanie domami i użytkownikami:} Możliwość tworzenia wirtualnego ,,domu'', zapraszania użytkowników poprzez kody lub linki oraz zarządzanie ich uprawnieniami.
    \item \textbf{System zadań (Tasks):} Tworzenie, przypisywanie i śledzenie statusu prac domowych.
    \item \textbf{Moduł finansowy (Bills):} Rejestrowanie rachunków, podział kosztów (równy lub według wag), śledzenie statusu płatności. Istotnym elementem jest integracja z mechanizmem OCR (Optical Character Recognition) pozwalającym na automatyczne wprowadzanie danych z fizycznych rachunków i faktur.
    \item \textbf{Listy zakupów (Shopping Lists):} Współdzielone listy produktów z aktualizacją w czasie rzeczywistym, co pozwala na dodawanie i odhaczanie pozycji przez wszystkich domowników.
    \item \textbf{System ankiet (Polls):} Narzędzie do demokratycznego podejmowania decyzji grupowych.
    \item \textbf{Zarządzanie inteligentnym domem (Smart Home):} Integracja, sterowanie oraz monitorowanie podłączonych urządzeń domowych z poziomu jednej aplikacji, co centralizuje kontrolę nad gospodarstwem domowym.
    \item \textbf{Powiadomienia:} System alertów przypominających o nadchodzących terminach płatności czy przydzielonych zadaniach.
\end{enumerate}

Aplikacja jest kierowana głównie do osób młodych, studentów wynajmujących wspólnie mieszkania oraz nowoczesnych rodzin, dla których smartfon jest naturalnym narzędziem organizacji życia.

\section{Analiza istniejących rozwiązań}

Na rynku dostępnych jest wiele narzędzi wspierających organizację, jednak większość z nich skupia się na wąskich dziedzinach lub jest adaptacją narzędzi biznesowych.

\subsection{Splitwise}
Splitwise to popularna aplikacja do dzielenia wydatków. Doskonale radzi sobie z matematyką podziału kosztów i upraszczaniem długów.
\begin{itemize}
    \item \textbf{Zalety:} Rozbudowane opcje podziału kosztów, obsługa wielu walut.
    \item \textbf{Wady:} Brak funkcji zarządzania zadaniami domowymi, listami zakupów czy komunikacji grupowej niezwiązanej z finansami.
\end{itemize}

\subsection{Trello / Asana}
Narzędzia typu Kanban, popularne w zarządzaniu projektami w IT.
\begin{itemize}
    \item \textbf{Zalety:} Elastyczność, możliwość tworzenia dowolnych tablic i list zadań.
    \item \textbf{Wady:} Zbyt skomplikowane do prostych zastosowań domowych, brak dedykowanych funkcji finansowych (np. podział rachunków), interfejs nastawiony na workflow biznesowy.
\end{itemize}

\subsection{OurGroceries}
Aplikacja dedykowana listom zakupów.
\begin{itemize}
    \item \textbf{Zalety:} Prosta obsługa, synchronizacja w czasie rzeczywistym.
    \item \textbf{Wady:} Ograniczona tylko do zakupów, przestarzały interfejs użytkownika (UI), brak integracji z finansami i zadaniami.
\end{itemize}

\subsection{Rozwiązania zintegrowane (np. OurHome, Cozi)}
Istnieją aplikacje próbujące łączyć te funkcje, jednak często borykają się z problemami:
\begin{itemize}
    \item Część z nich jest płatna w modelu subskrypcyjnym, co zniechęca studentów.
    \item Wiele projektów nie jest już rozwijanych i posiada przestarzały design.
    \item Rzadko oferują zaawansowane funkcje jak skanowanie rachunków (OCR).
\end{itemize}

\section{Podsumowanie analizy}

Analiza rynku wskazuje na brak darmowego, nowoczesnego i kompleksowego rozwiązania, które łączyłoby zarządzanie finansami (z wykorzystaniem OCR), zadaniami i zakupami w jednej aplikacji o wysokiej jakości interfejsu (UX/UI). Projektowany system ,,Household Manager'' ma na celu wypełnienie tej luki, oferując użytkownikom narzędzie typu ,,wszystko w jednym'', dostosowane do współczesnych standardów aplikacji mobilnych.

\section{Metodyka pracy i zarządzanie projektem}

W celu zapewnienia wysokiej efektywności prac oraz terminowej realizacji wszystkich założeń projektowych, proces wytwarzania oprogramowania był prowadzony zgodnie z metodyką zwinnego zarządzania projektami (ang. \textit{Agile}).

\subsection{Zastosowanie podejścia Kanban}

Do zarządzania zadaniami wybrano podejście \textbf{Kanban}, które charakteryzuje się wizualizacją przepływu pracy oraz skupieniem na ciągłym dostarczaniu kolejnych funkcjonalności. Wybór ten był podyktowany elastycznością metody, która pozwala na dynamiczne reagowanie na napotkane problemy techniczne oraz płynne przechodzenie między etapami implementacji warstwy serwerowej i klienckiej.

\subsection{Narzędzia wspomagające --- GitHub Projects}

Głównym narzędziem wykorzystywanym do zarządzania cyklem życia projektu była platforma \textbf{GitHub}, a w szczególności usługa \textbf{GitHub Projects}. Wszystkie wymagania funkcjonalne zostały zdekomponowane na mniejsze, atomowe jednostki pracy (ang. \textit{Issues}), które następnie trafiły do rejestru zadań produktu (ang. \textit{Backlog}).

Proces pracy został odzwierciedlony na cyfrowej tablicy Kanban składającej się z kolumn:
\begin{itemize}
    \item \textbf{Todo (Do zrobienia):} Zadania zaplanowane do realizacji, czekające w kolejce.
    \item \textbf{In Progress (W toku):} Funkcjonalności będące w danej chwili implementowane.
    \item \textbf{In Review (W weryfikacji):} Zadania ukończone programistycznie, przechodzące etap testów manualnych lub automatycznych.
    \item \textbf{Done (Ukończone):} W pełni przetestowane i wdrożone moduły aplikacji.
\end{itemize}

\begin{figure}[H]
    \centering
    \includegraphics[width=0.9\textwidth]{backlog.png}
    \caption{Tablica Kanban projektu w systemie GitHub Projects}
    \label{fig:backlog}
\end{figure}

Takie podejście umożliwiło precyzyjne śledzenie postępów prac oraz zapewniło transparentność całego procesu, co jest kluczowe przy realizacji złożonych systemów typu full-stack. Dzięki powiązaniu zadań z systemem kontroli wersji Git, każda zmiana w kodzie była bezpośrednio mapowana na konkretne wymaganie biznesowe zawarte w projekcie.